\chapter{Source conventions and the general Krawtchouk transform}
\label{chap:pq-source-summary}
This chapter integrates the supplied summary paper \emph{On Krawtchouk
Transforms} (\texttt{KravPQ.tex}) with \texttt{KravSums.tex}.
The originals are retained in \texttt{sources/}. Source conventions are
preferred here; explicit dictionaries keep the earlier symmetric chapters
compatible. Statements below are derived directly from the supplied
generating functions, rather than imported with an implicit normalization.

\section{The two source conventions}
In the transform paper, let $p,q>0$, $p+q=1$, and $x=N-2j$. Its convention is
\begin{equation}\label{eq:pq-source-gf}
G(v)=(1+2qv)^{N-j}(1-2pv)^j
=\sum_{n=0}^N\frac{v^n}{n!}K_n(x,N).
\end{equation}
In the sums paper the convention is
\begin{equation*}
(1+z)^{N-j}(1-rz)^j=\sum_{n=0}^Nz^n\phi^N_{nj},
\qquad
(1+z)^{N-j}(1-z)^j=\sum_{n=0}^Nz^n\Phi^N_{nj}.
\end{equation*}
Thus the lowercase $\phi$ denotes general parameters, and uppercase
$\Phi$ denotes the symmetric specialization in the sums convention.
The precise bridge, with the parameters of the transform paper, is
\begin{equation}\label{eq:pq-source-dictionary}
r=\frac pq,\qquad z=2qv,\qquad
K_n(N-2j,N)=n!(2q)^n\phi^N_{nj}.
\end{equation}
\emph{Parameter warning.} The sums paper writes $r=q_{\rm S}/p_{\rm S}$.
Consequently $p_{\rm S}=q$ and $q_{\rm S}=p$ when matching it to the
transform paper. Identifying the two letters $p$ without this interchange
would give an incorrect biased-parameter dictionary.
At $p=q=1/2$, $K_n(N-2j,N)=n!\Phi^N_{nj}$.
The earlier chapters' ordinary-coefficient notation $K_n(j;N)$ equals
$\Phi^N_{nj}$ in that symmetric setting; it is not the factorial-normalized
$K_n(x,N)$ of this chapter.

\section{Orthogonality and inversion}
Let $J$ have probabilities $\binom Nj q^jp^{N-j}$ and $X=N-2J$.
Independence of the $N$ Bernoulli steps gives
\begin{equation*}
\mathbb E[G(v)G(w)]
=\{p(1+2qv)(1+2qw)+q(1-2pv)(1-2pw)\}^N
=(1+4pqvw)^N.
\end{equation*}
Comparing coefficients proves
\begin{equation}
\mathbb E[K_n(X,N)K_m(X,N)]
=\delta_{nm}(n!)^2\binom Nn(4pq)^n.
\end{equation}
To distinguish ordinary from exponential coefficients, put
$\widehat\Phi_{nj}=K_n(N-2j,N)/n!$ and
\begin{equation*}
B=\operatorname{diag}\!\left(\binom Nj p^{N-j}q^j\right),\qquad
\Gamma=\operatorname{diag}\!\left(\binom Nn(4pq)^n\right).
\end{equation*}
These local source-style matrices are not the unweighted $B_N$ of the
symmetric chapters. Orthogonality and its inverse formula become
\begin{equation*}
\widehat\Phi B\widehat\Phi^{\mathsf T}=\Gamma,
\qquad
\widehat\Phi^{-1}=B\widehat\Phi^{\mathsf T}\Gamma^{-1}.
\end{equation*}
Indeed the first identity gives a right inverse of a square matrix.
In particular, this binomially weighted row orthogonality must not be
confused with the unweighted column sums in the sums chapter.

\section{Operator coefficients and positivity}
Write $D=d/dx$. For $\deg f\le N$ there is a unique expansion
\begin{equation*}
f(x)=\sum_{n=0}^N\widetilde f(n)K_n(x,N),\qquad
\widetilde f(n)=\frac1{n!}
\left[(pe^D+qe^{-D})^{N-n}(\sinh D)^nf\right]_{x=0}.
\end{equation*}
All power series act finitely on the polynomial space. To prove the formula,
put $M(z)=pe^z+qe^{-z}$ and $V(z)=\sinh z/M(z)$. Substitution in the generating
function gives
\begin{equation*}
e^{zx}=M(z)^N\sum_{n\ge0}\frac{V(z)^n}{n!}K_n(x,N).
\end{equation*}
For coefficient extraction through degree $N$, only $n\le N$ is needed;
compare coefficients of $z^m$ for $0\le m\le N$ to obtain the claimed
operator formula on monomials and hence on every such $f$.

\begin{proposition}
If $p\ge q$ and $f$ has nonnegative monomial coefficients, then
$\widetilde f(n)\ge0$ for every $n$. If $p>q$ and every monomial
coefficient of $f$ through degree $N$ is strictly positive, then every
$\widetilde f(n)$ is strictly positive.
\end{proposition}
\begin{proof}
The series $M(z)=\cosh z+(p-q)\sinh z$ and $\sinh z$ have nonnegative
coefficients. If $h(z)=\sum h_mz^m$ and $f(x)=\sum a_mx^m$, then
$[h(D)f]_{x=0}=\sum m!h_ma_m$. This proves nonnegativity.
For strict positivity the coefficient of $z^n$ in
$M(z)^{N-n}(\sinh z)^n$ is $1$, so the term $a_n$ already suffices.
\end{proof}
Strict positivity cannot be asserted for arbitrary polynomials with only
some positive coefficients: $f=1$ has zero transform coefficients for $n>0$.
At $p=q$, parity also forces zeros for many individual monomials.

\section{Low-degree checks and connection to sums}
Set $a=p-q$. Expansion of the generating function gives
\begin{equation*}
K_0=1,\qquad K_1=x-Na,\qquad
K_2=(x-Na)^2+2a(x-Na)-4Npq.
\end{equation*}
Consequently
\begin{equation*}
x=K_1+NaK_0,
\qquad
x^2=K_2+2a(N-1)K_1+(N^2a^2+4Npq)K_0.
\end{equation*}
These examples exhibit positivity when $p\ge q$ and recover
$x^2=K_2+NK_0$ in the symmetric case.

The weighted partial column sums, Catalan evaluations, and Boolean-lattice
applications of \texttt{KravSums.tex} are already developed in the sums
section included in the existing algebra chapter. Their entries should
be read as $\phi^N_{nj}$ (general) and $\Phi^N_{nj}$ (symmetric).
For example the bridge for any scalar weights $w_n$ is
\begin{equation*}
\sum_{n=0}^m w_n(\phi^N_{nj})^2
=\sum_{n=0}^m\frac{w_n}{(n!)^2(2q)^{2n}}
K_n(N-2j,N)^2.
\end{equation*}
This termwise conversion preserves the weights; replacing only the
polynomial symbol would not. It connects the column identities to the
operator transform without duplicating the existing sums treatment.

\section{Scope of this integration}
The supplied transform paper also treats coding theory, modified inversion,
and multivariate symmetric representations. This chapter integrates its
normalization, orthogonality, expansion calculus, positivity, and their
connection with the sums results. The remaining applications are retained
in the original reference source and are not claimed to be fully integrated
here. The malformed intermediate displays in the supplied Appell-inverse
calculation are not copied: the valid identity used above is
$V(z)=\sinh z/(pe^z+qe^{-z})$.
